\Needspace{7\baselineskip}
\section{The Intrinsic Three-Rank Spectrum}
\label{sec:12}
The construction has now reached the point at which a presentation-independent three-rank spectrum can be formed. Nothing in this section is selected by a charged-lepton measurement. The starting data are the rooted one-cell record, its two-active/one-idle coupling, the cyclic three-channel action, and the already established support coefficient

\begin{equation*}
\begin{adjustbox}{max width=\linewidth,center}
$\displaystyle
m=\frac{1}{2\pi}>0.
$
\end{adjustbox}
\end{equation*}
The lowercase symbol \(\lambda\) denotes the three intrinsic spectral amplitudes in this publication manuscript. Section 4 uses \(\xi\) for a lattice-site label. The labels attached to a particular displayed sector are not intrinsic; only the unordered spectrum and its minimum, middle, and maximum ranks survive a change of presentation.

Figure \(\ref{fig:c3-spectral-projections}\) depicts the three cyclic projections of one coefficient point.  It does not depict three independent points or the sides of a spatial triangle.

\Needspace{5\baselineskip}
\subsection{Rooted Action and the Spectral Phase}
Let

\begin{equation*}
\begin{adjustbox}{max width=\linewidth,center}
$\displaystyle
V=\mathbb C[\mathbb Z/3\mathbb Z]
$
\end{adjustbox}
\end{equation*}
with orthonormal basis \((e_0,e_1,e_2)\), and let \(P\) be the cyclic shift

\begin{equation*}
\begin{adjustbox}{max width=\linewidth,center}
$\displaystyle
Pe_r=e_{r+1\pmod 3},
\qquad
P^3=\mathbf1_V,
\qquad
P^\dagger=P^{-1}.
$
\end{adjustbox}
\end{equation*}
% LAYOUT_ONLY_GUARD:ROOTED_PROJECTOR_LEADIN
\Needspace{8\baselineskip}
The rooted two-active/one-idle projector is

\begin{equation*}
\begin{adjustbox}{max width=\linewidth,center}
$\displaystyle
\Pi_o
=
|0\rangle\langle0|
+
|1\rangle\langle1|
=
\begin{pmatrix}
1&0&0\\
0&1&0\\
0&0&0
\end{pmatrix}.
$
\end{adjustbox}
\end{equation*}
It satisfies

\begin{equation*}
\begin{adjustbox}{max width=\linewidth,center}
$\displaystyle
\Pi_o^\dagger=\Pi_o,
\qquad
\Pi_o^2=\Pi_o,
\qquad
\operatorname{rank}\Pi_o=2,
\qquad
\operatorname{Tr}\Pi_o=2.
$
\end{adjustbox}
\end{equation*}
% LAYOUT_ONLY_GUARD:CYCLIC_AVERAGING_PAIR
\Needspace{36\baselineskip}
Averaging its three cyclic presentations gives

\begin{equation*}
\begin{adjustbox}{max width=\linewidth,center}
$\displaystyle
\frac13
\sum_{r=0}^{2}
P^r\Pi_oP^{-r}
=
\frac23\mathbf1_V.
$
\end{adjustbox}
\end{equation*}
The normalized cyclic valuation independently gives

\begin{equation*}
\begin{adjustbox}{max width=\linewidth,center}
$\displaystyle
\nu(\Pi_o)=\frac23.
$
\end{adjustbox}
\end{equation*}
Trace preservation fixes the one-pass action aperture, and the completed unit read fixes the principal spectral phase:

\begin{equation}
\label{eq:action-phase}
\begin{adjustbox}{max width=\linewidth,center}
$\displaystyle
A_o^\#=\frac23,
\qquad
\Phi_o=A_o^\#I_{\mathrm{bulk},o}=\frac23,
\qquad
\theta_o:=\frac{\Phi_o}{3}=\frac29.
$
\end{adjustbox}
\end{equation}
Equation \(\eqref{eq:action-phase}\) is the phase source used below. As Section 10 established, this rooted projector/action derivation, not the separately displayed flow invariant, supplies that phase.

With the carried orientation \(\sigma\in\{+1,-1\}\), the action character is

\begin{equation*}
\begin{adjustbox}{max width=\linewidth,center}
$\displaystyle
\rho_{o,\sigma}
=
\exp\left(i\sigma\frac23\right).
$
\end{adjustbox}
\end{equation*}
Its three cube roots are

\begin{equation*}
\begin{adjustbox}{max width=\linewidth,center}
$\displaystyle
\zeta_{o,\sigma,j}
=
\exp\left[
i\left(
\frac{2\sigma}{9}
+
\frac{2\pi j}{3}
\right)
\right],
\qquad
j\in\mathbb Z/3\mathbb Z.
$
\end{adjustbox}
\end{equation*}
The factor \(\theta_o=2/9\) is therefore not fitted. It is one third of the already derived action phase.

\Needspace{5\baselineskip}
\subsection{Rooted Radius and Operator Coefficient}
The coefficient used in the spectral operator is fixed by the rooted two-active/one-idle coupling. For the escapement occurrence \(o\),

\begin{equation*}
\begin{adjustbox}{max width=\linewidth,center}
$\displaystyle
k_e(o)=Q_o(1,1,0)^{\mathsf T},
\qquad
Q_o^{\mathsf T}Q_o=\mathbf1_3,
\qquad
\|k_e(o)\|^2=2.
$
\end{adjustbox}
\end{equation*}
In the promoted response-chart notation, \(u,v\in\mathbb R\) are the response-plane coordinates and \(a\ne0\) is the chart-normalization scalar. Let \(T_o\) be the isometry from the active two-plane into the response plane. On that active plane,

\begin{equation*}
\begin{adjustbox}{max width=\linewidth,center}
$\displaystyle
z_{A(o)}
=
T_ok_e(o)
=
\left(
\frac{2u}{|a|},
\frac{2v}{|a|}
\right),
\qquad
T_o^{\mathsf T}T_o=\mathbf1_{\mathrm{active}},
\qquad
a\ne0.
$
\end{adjustbox}
\end{equation*}
Consequently,

\begin{equation}
\label{eq:rooted-radius}
\begin{adjustbox}{max width=\linewidth,center}
$\displaystyle
2
=
\|z_{A(o)}\|^2
=
\frac{4(u^2+v^2)}{|a|^2},
\qquad
\boxed{
u^2+v^2=\frac{|a|^2}{2}.
}
$
\end{adjustbox}
\end{equation}
This rooted-radius identity is recorded as Eq. \(\eqref{eq:rooted-radius}\).

Writing the permitted unit direction as \(\zeta_{o,\sigma,j}\), with \(|\zeta_{o,\sigma,j}|=1\), gives the exact assembly

\begin{equation*}
\begin{adjustbox}{max width=\linewidth,center}
$\displaystyle
u+iv
=
\frac{|a|}{\sqrt2}\,
\zeta_{o,\sigma,j}.
$
\end{adjustbox}
\end{equation*}
Only then is the positive ratio chart selected:

\begin{equation*}
\begin{adjustbox}{max width=\linewidth,center}
$\displaystyle
a=1
\quad\Longrightarrow\quad
\boxed{
u+iv=\frac{\zeta_{o,\sigma,j}}{\sqrt2}.
}
$
\end{adjustbox}
\end{equation*}
The \(1/\sqrt2\) coefficient below is therefore the normalized rooted radius, not a fitted parameter.

\Needspace{5\baselineskip}
\subsection{Six Presentations of One Intrinsic Object}
For each orientation and cubic-root branch, define the Hermitian operator

\begin{equation*}
\begin{adjustbox}{max width=\linewidth,center}
$\displaystyle
A_{\sigma,j}
=
\mathbf1_V
+
\frac1{\sqrt2}\zeta_{o,\sigma,j}P
+
\frac1{\sqrt2}\overline{\zeta_{o,\sigma,j}}P^{-1}.
$
\end{adjustbox}
\end{equation*}
% LAYOUT_ONLY_GUARD:FOURIER_DIAGONALIZATION
\Needspace{7\baselineskip}
Fourier diagonalization gives

\begin{equation*}
\begin{adjustbox}{max width=\linewidth,center}
$\displaystyle
\lambda^{\rm scr}_{\sigma,j;k}
=
1+\sqrt2\cos\left(
\sigma\theta_o
+
\frac{2\pi j}{3}
+
\frac{2\pi k}{3}
\right),
\qquad
k\in\mathbb Z/3\mathbb Z.
$
\end{adjustbox}
\end{equation*}
Changing \((\sigma,j)\) permutes these three values. The six permitted frames form the six-element \(D_3\) action, so the displayed operators are unitarily conjugate presentations of one orbit. No labeled sector is intrinsic.

The presentation-independent amplitude object is

\begin{equation}
\label{eq:intrinsic-spectrum}
\begin{adjustbox}{max width=\linewidth,center}
$\displaystyle
\operatorname{Spec}_{\mathrm{unord}}\mathfrak A_{\mathrm{int}}(o)
=
\left\{
1+\sqrt2\cos\left(
\theta_o+\frac{2\pi k}{3}
\right):
k\in\mathbb Z/3\mathbb Z
\right\}.
$
\end{adjustbox}
\end{equation}
Equation \(\eqref{eq:intrinsic-spectrum}\) is the intrinsic unordered spectrum; its displayed sector labels are not intrinsic.

For the reference calculation use the already fixed principal phase \(\theta_o=2/9\) and write

\begin{equation*}
\begin{adjustbox}{max width=\linewidth,center}
$\displaystyle
\lambda_0
=
1+\sqrt2\cos\theta_o,
$
\end{adjustbox}
\end{equation*}
\begin{equation*}
\begin{adjustbox}{max width=\linewidth,center}
$\displaystyle
\lambda_1
=
1-\frac{\cos\theta_o}{\sqrt2}
-\frac{\sqrt6}{2}\sin\theta_o,
$
\end{adjustbox}
\end{equation*}
\begin{equation*}
\begin{adjustbox}{max width=\linewidth,center}
$\displaystyle
\lambda_2
=
1-\frac{\cos\theta_o}{\sqrt2}
+\frac{\sqrt6}{2}\sin\theta_o.
$
\end{adjustbox}
\end{equation*}
Since \(0<\theta_o<\pi/12\),

\begin{equation*}
\begin{adjustbox}{max width=\linewidth,center}
$\displaystyle
\lambda_2-\lambda_1
=
\sqrt6\sin\theta_o
>0,
$
\end{adjustbox}
\end{equation*}
\begin{equation*}
\begin{adjustbox}{max width=\linewidth,center}
$\displaystyle
\lambda_0-\lambda_2
=
\sqrt{\frac32}
\left(
\sqrt3\cos\theta_o-\sin\theta_o
\right)
>0.
$
\end{adjustbox}
\end{equation*}
Moreover \(2\pi/3+\theta_o<3\pi/4\) implies

\begin{equation*}
\begin{adjustbox}{max width=\linewidth,center}
$\displaystyle
\lambda_1
>
1+\sqrt2\cos\frac{3\pi}{4}
=0.
$
\end{adjustbox}
\end{equation*}
The intrinsic rank values are consequently

\begin{equation*}
\begin{adjustbox}{max width=\linewidth,center}
$\displaystyle
\boxed{
0<
\lambda_{\min}
<
\lambda_{\mathrm{mid}}
<
\lambda_{\max},
}
$
\end{adjustbox}
\end{equation*}
with the reference-frame identification

\begin{equation*}
\begin{adjustbox}{max width=\linewidth,center}
$\displaystyle
(\lambda_{\min},\lambda_{\mathrm{mid}},\lambda_{\max})
=
(\lambda_1,\lambda_2,\lambda_0).
$
\end{adjustbox}
\end{equation*}
\Needspace{5\baselineskip}
\subsection{The Quadratic-Support Law}
Let

\begin{equation*}
\begin{adjustbox}{max width=\linewidth,center}
$\displaystyle
C=\{0,1,2\},
\qquad
J=C\times C,
$
\end{adjustbox}
\end{equation*}
and let \(E_{rp}=|r\rangle\langle p|\) be the canonical matrix units of the counting form. Their contraction law is

\begin{equation*}
\begin{adjustbox}{max width=\linewidth,center}
$\displaystyle
E_{rp}^\dagger E_{sq}
=
\delta_{rs}E_{pq}.
$
\end{adjustbox}
\end{equation*}
For \(A=\sum_{r,p\in C}A_{rp}E_{rp}\), the construction-derived adjoint and same-root contraction give

\begin{equation*}
\begin{adjustbox}{max width=\linewidth,center}
$\displaystyle
A^\dagger A
=
\sum_{p,q}
\left(
\sum_r\overline{A_{rp}}A_{rq}
\right)
E_{pq}.
$
\end{adjustbox}
\end{equation*}
The repeated root leg is evaluated while both position legs remain. A cross-root kernel would insert new structure, whereas tracing and redistributing would erase retained position information. The resulting map is

\begin{equation}
\label{eq:quadratic-support}
\begin{adjustbox}{max width=\linewidth,center}
$\displaystyle
\boxed{
M_{\mathrm{read}}(A)=mA^\dagger A,
\qquad
m=\frac1{2\pi}.
}
$
\end{adjustbox}
\end{equation}
The quadratic-support map in Eq. \(\eqref{eq:quadratic-support}\) is applied before any finite-screen rendering.

The complete nine-incidence book fixes the response power. With

\begin{equation*}
\begin{adjustbox}{max width=\linewidth,center}
$\displaystyle
U_J=\sum_{r,p\in C}E_{rp}
=\mathbf1_V+P+P^{-1}
=3\Pi_{\mathrm{sym}},
$
\end{adjustbox}
\end{equation*}
where \(\Pi_{\mathrm{sym}}\) is the publication alias for the projector onto the symmetric Fourier mode. This symbol is distinct from the rooted occurrence projector \(\Pi_o\). One has

\begin{equation*}
\begin{adjustbox}{max width=\linewidth,center}
$\displaystyle
\operatorname{Tr}(U_J^\dagger U_J)=9.
$
\end{adjustbox}
\end{equation*}
Within the positive-real family \(\mathcal F_s(B)=mB^s\), full-book closure requires

\begin{equation*}
\begin{adjustbox}{max width=\linewidth,center}
$\displaystyle
m9^s=9m.
$
\end{adjustbox}
\end{equation*}
Because \(m>0\), \(9>1\), and \(9^s\) is strictly increasing, this forces \(s=1\). No additional exponent is fitted.

Define the intrinsic quadratic supports

\begin{equation*}
\begin{adjustbox}{max width=\linewidth,center}
$\displaystyle
\boxed{
q_{\min}=m\lambda_{\min}^2,
\qquad
q_{\mathrm{mid}}=m\lambda_{\mathrm{mid}}^2,
\qquad
q_{\max}=m\lambda_{\max}^2.
}
$
\end{adjustbox}
\end{equation*}
Positivity preserves their strict rank. The common coefficient cancels from the exact ratios:

\begin{equation*}
\begin{adjustbox}{max width=\linewidth,center}
$\displaystyle
\boxed{
R_{\mathrm{mid}/\min}^{\mathrm{int}}
=
\left(
\frac{\lambda_{\mathrm{mid}}}{\lambda_{\min}}
\right)^2,
\qquad
R_{\max/\min}^{\mathrm{int}}
=
\left(
\frac{\lambda_{\max}}{\lambda_{\min}}
\right)^2.
}
$
\end{adjustbox}
\end{equation*}
These are exact, presentation-independent, dimensionless results. No particle name, measured mass, reporting unit, or finite-screen factor has entered their derivation.

